% =========================================================================================
% Appendices
% =========================================================================================
\appendix
\renewcommand{\thesubsection}{\Alph{subsection}}
\setcounter{subsection}{0}

\section*{Appendices}
\addcontentsline{toc}{section}{Appendices}

\subsection{Dataset Details}
\label{app:datasets}

Chapter~\ref{sec:data} gives composition and provenance. This appendix records the detail
needed to reproduce the partitions.

\paragraph{Exact counts after cleaning.}

\begin{table}[H]
  \centering
  \small
  \begin{tabularx}{\textwidth}{@{}lXrr@{}}
    \toprule
    Dataset & Unit recovered & Files & Cases \\
    \midrule
    Lung (POCUS-style clips) & clip, then case by filename pattern & 187 clips & 165 \\
    Gallbladder (teaching atlas) & case by directory, capped at 25 frames & 4\,826 frames & 199 \\
    Cardiac (CAMUS) & patient, provided by the dataset & --- & 100 (validation) \\
    Triage (pooled) & encounter & --- & 16\,180 (test) \\
    \bottomrule
  \end{tabularx}
  \caption{Counts behind the reported figures. Frames are capped per gallbladder case so that
    longer studies do not dominate the loss.}
  \label{tab:app-counts}
\end{table}

\paragraph{Recovering case identity for the lung dataset.} The release carries no patient
identifier linking clips belonging to one patient, so case identity was reconstructed from the
filename pattern, which encodes the source acquisition. Clips sharing that prefix are treated
as one case and never separated across folds. This is a heuristic: it can merge two patients
whose files happen to share a prefix, or split one patient recorded under two. The residual
uncertainty is stated as a limitation in Section~\ref{subsec:results-limitations} and is the
reason the lung case count is reported as approximate rather than exact.

\paragraph{Class definitions.} The gallbladder taxonomy trains on eight classes ---
gallstones, cholecystitis, membranous/gangrenous, perforation, polyps, adenomyomatosis,
carcinoma, wall thickening --- and reports five by summing group probabilities:
cholelithiasis; acute cholecystitis (any severity), formed from cholecystitis, membranous and
perforation; polyps and adenomyomatosis; carcinoma; and wall thickening. A ninth folder present
in the original release was excluded after the per-case inspection described in
Section~\ref{subsec:gallbladder-module}.

The lung module models four findings: B-lines, consolidation, pleural effusion and pleural
thickening. Pneumothorax is not modelled --- the dataset contains no positive cases.

\paragraph{Licences.} Conditions of use are stated per dataset in Section~\ref{subsec:data-overview}.

\subsection{Model Architectures}
\label{app:architectures}

\begin{table}[H]
  \centering
  \small
  \begin{tabularx}{\textwidth}{@{}lXr@{}}
    \toprule
    Module & Architecture & Parameters \\
    \midrule
    Cardiac & U-Net, EfficientNet-B0 encoder, 4 output classes & 6\,251\,328 \\
    Gallbladder & EfficientNet-B0, 8-way classifier head & 4\,017\,796 \\
    Lung & EfficientNet-B0 features $\rightarrow$ dropout(0.4) $\rightarrow$ 512 ReLU
      $\rightarrow$ dropout(0.4) $\rightarrow$ 4 & 4\,665\,472 \\
    Triage & Gradient-boosted trees (XGBoost) & --- \\
    \bottomrule
  \end{tabularx}
  \caption{Architectures and parameter counts. The lung head is deliberately wider than a
    linear classifier because the backbone's last blocks are unfrozen.}
  \label{tab:app-arch}
\end{table}

\paragraph{Training hyperparameters.} All three imaging modules train with mixed precision, an
AdamW optimiser and a cosine learning-rate schedule. Input size is $224\times224$ for the two
classifiers and $256\times256$ for segmentation.

\begin{table}[H]
  \centering
  \small
  \begin{tabular}{@{}llll@{}}
    \toprule
    & Lung & Cardiac & Gallbladder \\
    \midrule
    Split & 3-fold GroupKFold by case & patient-wise & 5-fold GroupKFold by case \\
    Epoch budget & 15 & --- & --- \\
    Early stopping & val AUROC, patience 3 & val Dice, best kept & --- \\
    Loss & masked BCE, multi-label & cross-entropy + soft Dice & weighted cross-entropy \\
    Calibration & Platt, per finding & isotonic & temperature ($T=0.87$) \\
    \bottomrule
  \end{tabular}
  \caption{Training configuration per imaging module. Class weights for the gallbladder module
    are computed per fold from the training portion only.}
  \label{tab:app-hyper}
\end{table}

The triage model uses a temporal split --- fit on the earlier year, test on the later --- with
class weighting favouring the urgent tiers, and a safety threshold of 0.20 on the high-urgency
probability.

\subsection{Additional Results}
\label{app:results}

\paragraph{Lung, per fold.} The pooled out-of-fold figures in Table~\ref{tab:lung-results} are
computed over predictions from all three folds together. The per-fold AUROC below shows how
much they vary, which is the honest indication of how stable a 165-case estimate is.

\begin{table}[H]
  \centering
  \small
  \begin{tabular}{@{}lcccc@{}}
    \toprule
    Finding & Fold 0 & Fold 1 & Fold 2 & Pooled \\
    \midrule
    B-lines & 0.805 & 0.536 & 0.551 & 0.830 \\
    Consolidation & 0.808 & 0.867 & 0.817 & 0.832 \\
    Pleural effusion & 0.745 & 0.593 & 0.738 & 0.718 \\
    Pleural thickening & 0.795 & 0.728 & 0.740 & 0.773 \\
    \bottomrule
  \end{tabular}
  \caption{Lung AUROC per fold against the pooled out-of-fold figure. B-lines varies from 0.536
    to 0.805 across folds --- a spread that a single pooled number conceals entirely, and the
    clearest available evidence that 165 cases is a small sample.}
  \label{tab:app-lung-folds}
\end{table}

The pooled figure is not the mean of the three folds. Pooling the out-of-fold predictions and
computing one AUROC over all cases is the correct aggregation; averaging per-fold AUROCs weights
small folds equally with large ones and is reported here only to expose the variance.

\paragraph{Triage, per tier.}

\begin{table}[H]
  \centering
  \small
  \begin{tabular}{@{}lcccr@{}}
    \toprule
    Tier & Precision & Recall & F1 & Support \\
    \midrule
    Low & 0.557 & 0.638 & 0.594 & 3\,452 \\
    Medium & 0.647 & 0.595 & 0.620 & 6\,263 \\
    High & 0.783 & 0.784 & 0.784 & 6\,465 \\
    \midrule
    Macro & 0.662 & 0.672 & 0.666 & 16\,180 \\
    \bottomrule
  \end{tabular}
  \caption{Triage performance per tier on the held-out later year.}
  \label{tab:app-triage}
\end{table}

\paragraph{Cardiac, ejection fraction bands.} Per-band recall on 100 validation patients:
severe 0.727 ($n=11$), reduced 0.697 ($n=66$), normal 0.478 ($n=23$). Severe-to-normal
confusions: 0 of 100.

\paragraph{Reliability diagrams} accompany each calibration in the module notebooks listed in
Appendix~\ref{app:implementation}. Expected calibration errors before and after are reported
in each module's calibration section in Chapter~\ref{sec:perception}.

\subsection{Confusion Matrices}
\label{app:confusion}

\paragraph{Gallbladder, five reported classes.} Per-class recall is given in
Table~\ref{tab:gb-results}. The structure of the errors matters more than the diagonal:
\textbf{18\% of the fine-grained model's confusions stay inside their reported group}, meaning
82\% cross a group boundary. That figure is what decided against reporting the coarser
three-category grouping as the module's output --- if the groups captured a distinction the
model could see, confusions would concentrate within them, and they do not.

\paragraph{The collapsed matrix.} The taxonomy selection of
Section~\ref{subsec:gallbladder-module} operated on the fine-grained confusion matrix directly:
for each candidate grouping, the matrix rows and columns belonging to one group were summed to
give the confusion matrix that grouping would produce, and balanced accuracy was read off
without training anything. The chosen candidate was predicted at 63.0\% and measured at 59.7\%
after training.

\paragraph{Triage.} The three-tier confusion is summarised by
Table~\ref{tab:app-triage}. The error that matters is asymmetric and is reported separately in
Section~\ref{subsec:triage-error}: dangerous under-triage, a high-acuity case predicted low,
occurred at 3.5\% without the safety threshold and 1.5\% with it.

\subsection{Example Grad-CAM Maps}
\label{app:gradcam}

Attention maps are produced by each imaging module and the aggregate statistics are reported in
Section~\ref{subsec:explainability}.

The selection principle for any figures placed here matters more than the figures themselves: a
selection of successful maps is not evidence, because a viewer who knows where the pathology
lies will read the map as pointing at it. Any maps reproduced should therefore include cases the
module got wrong, so that a reader can see whether the attention differs between correct and
incorrect predictions --- which is the only comparison from which anything can be concluded
without expert annotation.

The measured statistics stand on their own regardless: centroid spread, area above
half-maximum and normalised entropy were computed over the full sets rather than over a chosen
subset, together with an explicit test for centre bias.

\subsection{Technical Implementation}
\label{app:implementation}

\paragraph{Repository structure.}

\begin{verbatim}
src/agents/
  schema.py              the contract all agents write through
  triage/                Triage Agent
  ultrasound/            Ultrasound Agent
    agent.py               organ routing; lung inference, importable and tested
  clinical/              Clinical Reasoning Agent
    clinical_state.py      structured state, evidence enumeration
    reasoning.py           escalation policy, prompt, validators
    llm.py                 llama.cpp backend plus test backends
    retrieval.py           TF-IDF retrieval with a relevance floor
    decision_support.py    severity, alerts, examinations, scenario routing
    report.py              standardised report and archiving
    thresholds.json        every alert cutoff, versioned
    corpus/                32 sourced knowledge units
    run_case.py            five benchmark scenarios, end to end
  tests/                 228 tests, grouped by safety property
tools/run_regression.py  150 end-to-end checks over 10 cases
app.py                   clinician-facing Streamlit application
notebooks/               training, one per organ
models/                  metrics, calibrators, frozen baselines
manifests/               per-source dataset manifests
\end{verbatim}

\paragraph{Environment.} Python 3.13, PyTorch 2.13, torchvision 0.28, scikit-learn 1.5,
XGBoost 3.2, NumPy 2.3. Segmentation uses \texttt{segmentation\_models\_pytorch}; the language
model runs through \texttt{llama-cpp-python} with a CUDA build. Training and inference ran on
Google Colab T4 GPUs.

\paragraph{Reproducing the training runs.} Each module has one notebook under
\texttt{notebooks/}, which loads its dataset, performs its own grouped split, trains, evaluates
and exports both weights and calibrator. The notebooks are self-contained and were the interface
through which every reported figure was produced.

\paragraph{Reproducing the safety benchmark.} The 228 tests require no GPU, no model weights and
no dataset:

\begin{verbatim}
python -m src.agents.tests.run_benchmark
\end{verbatim}

Test modules are discovered from the directory rather than listed, after a hand-maintained list
silently omitted a new file of fourteen tests and the suite reported the same total as before,
all passing.

\paragraph{Reproducing the reasoning evaluation.} \texttt{run\_case} executes one scenario end
to end. Its \texttt{-{}-dry-run} mode prints the clinical state, the escalation decision and the
exact prompt without loading the language model:

\begin{verbatim}
python -m src.agents.clinical.run_case --scenario conflict --dry-run
python -m src.agents.clinical.run_case --scenario conflict --n-gpu-layers -1
\end{verbatim}

Determinism is enforced by temperature 0, a fixed seed and a KV-cache reset before each call;
two identical calls produce byte-identical output.

\paragraph{The output schema.} A report carries the organ, a status, a list of detected findings
with calibrated confidences, a list of findings screened for and not detected, and a reliability
block declaring scope. Section~\ref{subsec:output-contract} states the four validation rules and
the reason for each.
